\documentclass[10pt,a4paper]{amsart}
\usepackage[margin=19mm]{geometry}
\usepackage{amsmath,amssymb,mathtools}
\usepackage[scr=rsfs,cal=euler]{mathalfa}
\usepackage{xfrac}
\usepackage[unicode,hidelinks,bookmarks=false]{hyperref}
\newcommand{\ie}{i.e.\ }
\newcommand{\cf}{cf.\ }
\newcommand{\resp}{respectively\ }
\newcommand{\Subsection}{\subsection}
\providecommand{\nobreakdash}{\nobreak\hbox{-}}
\setlength{\emergencystretch}{2em}
\multlinegap0pt
\usepackage{bm}
\usepackage{bbm}
\usepackage{dsfont}
\usepackage{xspace}
\usepackage{cancel}
\usepackage{mathtools}
\usepackage{amsthm}
\makeatletter
\@ifundefined{theorem}{\newtheorem{theorem}{Theorem}}{}
\@ifundefined{lemma}{\newtheorem{lemma}[theorem]{Lemma}}{}
\makeatother
\title[Active-set Newton-MR methods for nonconvex optimization prob]{Active-set Newton-MR methods for nonconvex optimization problems with bound constraints}
\author{Ernesto G. Birgin, Geovani N. Grapiglia, Diaulas S. Marcondes}
\date{Source: arXiv:2508.20967v1 (CC BY 4.0)}
\begin{document}
\maketitle

\noindent\emph{Source and licence.} Active-set Newton-MR methods for nonconvex optimization problems with bound constraints. Version arXiv:2508.20967v1 (CC BY 4.0), distributed under the Creative Commons Attribution 4.0 International licence, as recorded in the arXiv licence metadata for this exact version and shown on the abstract page https://arxiv.org/abs/2508.20967v1. Full rights notice: licenses/arxiv-2508.20967-CC-BY-4.0.txt.\par\smallskip

\noindent\textit{Editorial scope.} Counted unit: the Lemma whose source label is \texttt{eq:edg1} in the article (source0.tex lines 740--759, source label \texttt{eq:edg1}), delivered together with the article's own complete original proof of it (source0.tex lines 761--895, one \texttt{proof} environment of 135 lines). No mathematics is written, repaired or re-derived: statement and proof are the article's own text line for line. Required context retained uncounted: eq:fx (lines 168--171); eq:2.3 (lines 232--234); eq:minres1 (lines 473--478); eq:minres3 (lines 473--478); eq:minres4 (lines 473--478); eq:flag (lines 520--531); eq:2.21 (lines 552--554). Every definition, display and label of those lines is retained unchanged. Omitted from this package and not redistributed: 1--167, 172--231, 235--472, 479--519, 532--551, 555--739, 760--760, 896--1619. Nothing inside a retained excerpt is omitted or altered: every retained body is delivered exactly as the article's lines are written, so no omission receipt of any kind accompanies this package. Citations: the retained excerpts carry no citation command at all, so no bibliography entry of the article is transcribed and no reference is redistributed. Reference adaptation: none was needed and none was made; every reference inside the delivered unit resolves inside it, so no unresolved reference or citation is produced. Printed slips kept literal: no printed word, symbol, hypothesis or number of the counted statement or of its proof was altered. Layout and class: the article's own class is \texttt{article} with packages including geometry, color, float, textcomp, url, amsthm, amsmath, amssymb, mdframed, multirow, changes, hyperref, mathtools, booktabs; the wrapper is the lane's pinned \texttt{amsart} editorial wrapper at 10pt on A4, which restates the theorem-like environments the retained text needs and adds the article's own macro definitions that the retained text needs (none), with extra wrapper packages (bm, bbm, dsfont, xspace, cancel, mathtools). The delivered numbering is the unit's own. Assets: no retained span contains a figure environment, a graphics-inclusion command, a PDF-page inclusion, a table or a TikZ picture; no image file is copied into this package, the package declares no asset dependency and no image file is redistributed with it. Licence: version arXiv:2508.20967v1 of this article is distributed under the Creative Commons Attribution 4.0 International licence, as recorded in the arXiv licence metadata for this exact version and shown on the abstract page https://arxiv.org/abs/2508.20967v1. A full rights notice is delivered at licenses/arxiv-2508.20967-CC-BY-4.0.txt. Characters: every retained excerpt is pure ASCII, so the delivered engine, XeLaTeX with its default Latin Modern text font, reports no missing character.
\begin{equation}
f_{x}(y)=f(x+Q(x)y).
\label{eq:fx}
\end{equation}

\begin{equation} \label{eq:2.3}
f_{x^k}((0.5)^{\ell_{k}} \alpha_0 d^k) \leq f_{x^{k}}(0)+\rho (0.5)^{\ell_{k}} \alpha_0 \langle\nabla f_{x^{k}}(0),d^{k}\rangle.
\end{equation}

\begin{eqnarray} 
&\langle g,s\rangle+\langle Hs,s\rangle\leq 0,& \label{eq:minres1} \\
&\|Hr\|\leq\eta\|Hs\|,& \label{eq:minres3} \\
&\langle g,r\rangle=-\|r\|^{2},& \label{eq:minres4} \\
&\langle Hr,r\rangle<0.& \label{eq:minres5}
\end{eqnarray}

\begin{equation} \label{eq:flag}
\texttt{Flag}=
\left\{
\begin{array}{ll} 
1, & \text{if } \texttt{D}_{\texttt{type}}^{k}=\texttt{`SOL'} 
\text{ and } \min\left\{\frac{\langle\nabla^{2}f_{x^{k}}(0)d_{1}^{k},d_{1}^{k}\rangle}{\|d_{1}^{k}\|^{2}},\frac{\langle\nabla^{2}f_{x^{k}}(0)r^{k},r^{k}\rangle}{\|r^{k}\|^{2}}\right\}<\eta_k,\\
1, & \text{if } \texttt{D}_{\texttt{type}}^{k}=\texttt{`NPC'} 
\text{ and } \|\nabla^{2}f_{x^{k}}(0)r^{k}\|\leq\eta_k\|\nabla^{2}f_{x^{k}}(0)s^{k}\|,\\
0, & \text{otherwise}.
\end{array}
\right.
\end{equation}

\begin{equation} \label{eq:2.21}
f_{x^k}((0.5)^{\ell_{k}} \alpha_0 d^k) \leq f_{x^{k}}(0)+\rho (0.5)^{\ell_{k}} \alpha_0 \langle\nabla f_{x^{k}}(0),d^{k}\rangle.
\end{equation}

\begin{lemma}
Suppose that A1 and A3 hold. Then, whenever $x^{k+1}$ is computed at Step 4 of Algorithm~5, we have
\small
\setlength{\arraycolsep}{3pt}
\begin{equation}
f(x^{k})-f(x^{k+1})\geq\left\{\begin{array}{ll}
\frac{\rho(1-\rho)}{L_{g}}\left(\frac{a_{2}}{a_{1}}\right)^{2}\|\nabla_{\Omega}^{I}f(x^{k})\|^{2},&\text{if $\texttt{Flag}=1$,}\\
& \\
\frac{\rho}{2^{5/2}}\left(\frac{\eta}{L_{g}}\right)^{3/2}\sqrt{\frac{3(1-2\rho)}{L_{H}}}\|\nabla_{\Omega}^{I}f(x^{k})\|^{3/2},&\text{if $\texttt{Flag}=0$ and $\texttt{D}_{\texttt{type}}^{k}=\texttt{`SOL'}$,}\\
&\\
\frac{\rho}{2\left[\left(\frac{L_{g}}{\eta}\right)+1\right]^{3/2}}\sqrt{\frac{6(1-\rho)}{L_{H}}}\|\nabla_{\Omega}^{I}f(x^{k})\|^{3/2},&\text{if $\texttt{Flag}=0$ and $\texttt{D}_{\texttt{type}}^{k}=\texttt{`NPC'}$.}
\end{array}
\right.
\label{eq:edg1}
\end{equation}
\normalsize
Moreover, the number of evaluations of $f(\,\cdot\,)$ required to guarantee the fulfilment of (\ref{eq:2.21}) is bounded from above by 
$$|\log_{2}\left(\min\left\{(1-\rho)\eta / 2L_{g},(1-\rho)a_{2} / 2L_{g}a_{1}^{2}\right\}\right)|.$$
\label{lem:edg1}
\end{lemma}

\begin{proof}
Since $x^{k+1}$ is computed by Step 4 of Algorithm 5, it follows that $\ell_{k}>0$, and so
\begin{equation}
f_{x^{k}}((0.5)^{\ell_{k}-1}\alpha_{0}d^{k})>f_{x^{k}}(0)+\rho(0.5)^{\ell_{k}-1}\alpha_{0}\langle\nabla f_{x^{k}}(0),d^{k}\rangle.
\label{eq:edg2}
\end{equation}
Let us analyse separately the possible cases.
\\[0.2cm]
\noindent\textbf{Case 1:} $\texttt{Flag}=1$.
\\[0.2cm]
\noindent By Step 1 of Algorithm 5, $d^{k}$ is defined by the same correction procedure employed in Algorithm 2. Consequently, by Lemma 2.1 we have
\begin{equation}
\|d^{k}\|\leq a_{1}\|\nabla f_{x^{k}}(0)\|\quad\text{and}\quad \langle\nabla f_{x^{k}}(0),d^{k}\rangle\leq -a_{2}\|\nabla f_{x^{k}}(0)\|^{2}.
\label{eq:edg3}
\end{equation}
Then, by A1, Lemma 2.2 also applies to $d^{k}$. In particular, it follows from (\ref{eq:edg2}) that
\begin{equation}
(0.5)^{\ell_{k}-1}\alpha_{0}> 2(1-\rho)a_{2} / (L_{g}a_{1}^{2}).
\label{eq:edg4}
\end{equation}
Now, combining (\ref{eq:fx}), (\ref{eq:2.21}), the second inequality in (\ref{eq:edg3}), and (\ref{eq:edg4}), we obtain
\small
\begin{eqnarray*}
f(x^{k})-f(x^{k+1})&=& f_{x^{k}}(0)-f_{x^{k}}((0.5)^{\ell_{k}}\alpha_{0}d^{k})\geq\rho (0.5)^{\ell_{k}}\alpha_{0}\left(-\langle\nabla f_{x^{k}}(0),d^{k}\rangle\right)\\
&\geq &\frac{\rho}{2}(0.5)^{\ell_{k}-1}\alpha_{0}a_{2}\|\nabla f_{x^{k}}(0)\|^{2}\geq\frac{\rho(1-\rho)}{L_{g}}\left(\frac{a_{2}}{a_{1}}\right)^{2}\|\nabla f_{x^{k}}(0)\|^{2}\\
&\geq &\frac{\rho (1-\rho)}{L_{g}}\left(\frac{a_{2}}{a_{1}}\right)^{2}\|\nabla_{\Omega}^{I}f(x^{k})\|^{2},
\end{eqnarray*}
\normalsize
that is, (\ref{eq:edg1}) holds in this case.
\\[0.2cm]
\noindent\textbf{Case 2:} $\texttt{Flag}=0$ and $\texttt{D}_{\texttt{type}}^{k}=\texttt{`SOL'}$.
\\[0.2cm]
\noindent In this case, it follows from (\ref{eq:edg2}) and Lemma 2.7 that
\small
\begin{equation}
(0.5)^{\ell_{k}-1}\alpha_{0}>\max\left\{\sqrt{3(1-2\rho)\eta/(L_{H}\|d^{k})\|},2(1-\rho)\eta/L_{g}\right\}.
\label{eq:edg5}
\end{equation}
\normalsize
In addition, by (\ref{eq:flag}), (\ref{eq:minres1})-(\ref{eq:minres3}), and $\eta_{k}\geq\eta$, we also have 
\begin{eqnarray}
\langle\nabla f_{x^{k}}(0),d^{k}\rangle&\leq &-\langle\nabla^{2}f_{x^{k}}(0)d^{k},d^{k}\rangle,
\label{eq:edg6}\\
\langle\nabla^{2}f_{x^{k}}(0)d^{k},d^{k}\rangle&\geq&\eta\|d^{k}\|^{2},
\label{eq:edg7}\\
\langle\nabla^{2}f_{x^{k}}(0)r^{k},r^{k}\rangle&\geq&\eta\|r^{k}\|^{2},
\label{eq:edg8}\\
\|\nabla^{2}f_{x^{k}}(0)r^{k}\|&\leq&\eta_{k}\|\nabla^{2}f_{x^{k}}(0)d^{k}\|.
\label{eq:edg9}
\end{eqnarray}
Combining (\ref{eq:fx}), (\ref{eq:2.21}), (\ref{eq:edg6}) and (\ref{eq:edg7}), it follows that
\begin{eqnarray}
    f(x^{k})-f(x^{k+1})&=&f_{x^{k}}(0)-f_{x^{k}}((0.5)^{\ell_{k}}\alpha_{0}d^{k})\geq \rho (0.5)^{\ell_{k}}\alpha_{0}\left(-\langle\nabla f_{x^{k}}(0),d^{k}\rangle\right)\nonumber\\
    &\geq &\frac{\rho}{2}(0.5)^{\ell_{k}-1}\alpha_{0}\langle\nabla^{2}f_{x^{k}}(0)d^{k},d^{k}\rangle\geq \frac{\eta\rho}{2}\sqrt{\frac{3(1-2\rho)\eta}{L_{H}\|d^{k}\|}}\|d^{k}\|^{2}.
\label{eq:edg10}
\end{eqnarray}
On the other hand, by (\ref{eq:edg8}) we have
\[
\|r^{k}\|^{2}=\frac{\|r^{k}\|^{2}}{\langle\nabla^{2}f_{x^{k}}(0),r^{k}\rangle}\langle\nabla^{2}f_{x^{k}}(0)r^{k},r^{k}\rangle\leq\frac{\|\nabla^{2}f_{x^{k}}(0)r^{k}\|\|r^{k}\|}{\eta_{k}}.
\]
Then, dividing both sides by $\|r^{k}\|$ and using (\ref{eq:edg9}) and A1 we get
\[
\|r^{k}\|\leq\frac{\|\nabla^{2}f_{x^{k}}(0)r^{k}\|}{\eta_{k}}\leq\|\nabla^{2}f_{x^{k}}(0)d^{k}\|\leq\|\nabla^{2}f_{x^{k}}(0)\|\|d^{k}\|\leq L_{g}\|d^{k}\|.
\]
Thus,
\begin{eqnarray}
\|\nabla f_{x^{k}}(0)\|&\leq&
\|r^{k}\|+\|\nabla^{2}f_{x^{k}}(0)d^{k}\|\leq 2L_{g}\|d^{k}\|.
\label{eq:edg11}
\end{eqnarray}
Combining (\ref{eq:edg10}) and (\ref{eq:edg11}), it follows that
\small
\begin{eqnarray*}
    f(x^{k})-f(x^{k+1})&\geq &\frac{\eta\rho}{2}\sqrt{\frac{3(1-2\rho)\eta}{L_{H}\|d^{k}\|}}\frac{\|\nabla f_{x^{k}}(0)\|^{3/2}}{(2L_{g})^{3/2}}\\
    &\geq & \frac{\rho}{2^{5/2}}\left(\frac{\eta}{L_{g}}\right)^{3/2}\sqrt{\frac{3(1-2\rho)}{L_{H}}}\|\nabla_{\Omega}^{I}f(x^{k})\|^{3/2},
\end{eqnarray*}
\normalsize
that is, (\ref{eq:edg1}) holds in this case.
\\[0.2cm]
\noindent\textbf{Case 3:} $\texttt{Flag}=0$ and $\texttt{D}_{\texttt{type}}^{k}=\texttt{`NPC'}$.
\\[0.2cm]
\noindent In this case, it follows from (\ref{eq:edg2}) and Lemma 2.8 that
\begin{equation}
(0.5)^{\ell_{k}-1}\alpha_{0}>\max\left\{\sqrt{6(1-\rho)/(L_{H}\|d^{k}\|)},2(1-\rho)/L_{g}\right\}.
\label{eq:edg12}
\end{equation}
In addition, by (\ref{eq:flag}) and (\ref{eq:minres4}) we have
\begin{equation}
    \langle\nabla f_{x^{k}}(0),d^{k}\rangle=-\|d^{k}\|^{2},
    \label{eq:edg13}
\end{equation}
and
\begin{equation}
\|\nabla^{2}f_{x^{k}}(0)d^{k}\|>\eta_{k}\|\nabla^{2}f_{x^{k}}(0)s^{k}\|\geq \eta\|\nabla^{2}f_{x^{k}}(0)s^{k}\| .
    \label{eq:edg14}
\end{equation}
Combining (\ref{eq:fx}), (\ref{eq:2.3}), (\ref{eq:edg13}) and (\ref{eq:edg12}), it follows that
\small
\begin{eqnarray}
f(x^{k})-f(x^{k+1})&\geq& 
\rho (0.5)^{\ell_{k}}\alpha\left(-\langle\nabla f_{x^{k}}(0),d^{k}\rangle\right)=\frac{\rho}{2}(0.5)^{\ell_{k}-1}\alpha_{0}\|d^{k}\|^{2}\nonumber\\
&\geq &(\rho/2)\sqrt{6(1-\rho)/L_{H}}\|d^{k}\|^{3/2}.
\label{eq:edg15}
\end{eqnarray}
\normalsize
On the other hand, by (\ref{eq:edg14}) and A1, we have
\small
\begin{equation} \label{eq:edg16}
\begin{array}{lclcl}
\|\nabla f_{x^{k}}(0)\| & \leq & \|\nabla f_{x^{k}}(0)+r^{k}\|+\|r^{k}\| & = & \|\nabla^{2} f_{x^{k}}(0)s^{k}\|+\|d^{k}\| \\[2mm]
& < & \|\nabla^{2} f_{x^{k}}(0)d^{k}\| / \eta+\|d^{k}\| & \leq & \|\nabla^{2}f_{x^{k}}(0)\|\|d^{k}\| / \eta +\|d^{k}\| \\[2mm]
& \leq & (L_{g}/\eta+1) \|d^{k}\|. & & 
\end{array}
\end{equation}
\normalsize
Combining (\ref{eq:edg15}) and (\ref{eq:edg16}), it follows that
\small
\[
f(x^{k})-f(x^{k+1}) 
\geq 
\frac{\rho}{2}\sqrt{\frac{6(1-\rho)}{L_{H}}}\frac{\|\nabla f_{x^{k}}(0)\|^{3/2}}{(L_{g}/\eta+1)^{3/2}}
\geq \frac{\rho}{2(L_{g}/\eta+1)^{3/2}}\sqrt{\frac{6(1-\rho)}{L_{H}}}\|\nabla_{\Omega}^{I}f(x^{k})\|^{3/2},
\]
\normalsize
that is, (\ref{eq:edg1}) also holds in this case.
\\[0.2cm]
To conclude, notice that the number of evaluations of $f(\,\cdot\,)$ performed at Step 3 of Algorithm 5 is equal to $\ell_{k}+1$. Since $\alpha_{0}\in (0,1]$, it follows from (\ref{eq:edg3}), (\ref{eq:edg5}) and (\ref{eq:edg12}) that
\[
(0.5)^{\ell_{k}+1}\geq (0.5)^{\ell_{k}+1}\alpha_{0}>\min\left\{(1-\rho)\eta/(2L_{g}), (1-\rho)a_{2}/(2L_{g}a_{1}^{2}),(1-\rho)/(2L_{g})\right\}.
\]
Then, taking the logarithm on both sides and using the fact that $\eta\in (0,1]$, we obtain
\[
\ell_{k}+1\leq\left|\log_{2}\left(\min\left\{(1-\rho)\eta/(2L_{g}),(1-\rho)a_{2}/(2L_{g}a_{1}^{2})\right\}\right)\right|.
\]
\end{proof}

\end{document}
